% !TeX program = xelatex
% UTF-8. Compile from this directory with XeLaTeX.
% Define \GDMInNotebook before input in an existing XePersian document.
% Source: Google Drive notebook/GDM, handwritten pages 1--5.
% Source statements and numbers are transcribed, not clinically revised.
% pic01.jpg: 20261016_212621014.jpg; 14vzYaNQewR-Zv15xJTy_B3uFVi6Z8FJj
% pic02.jpg: 20261016_212643522.jpg; 1wjIUmEzotv3aRGP_wFFXDe7rv3BFZXa3 (handwritten page 2)
% pic03.jpg: 20261016_212635005.jpg; 1vEpxo0twwy0h30pxoKh3Ywu5VjlXJXrZ (handwritten page 3)
% pic04.jpg: 20261016_212643522.jpg; 1lbWWO_zL8ZRKzIA3KOSICeCS3S7x4I00
% pic05.jpg: 20261016_212649498.jpg; 1Scv6hEVBvyXbnsyDiWGeLic0pgC_8xBo
\ifdefined\GDMInNotebook
  \let\GDMfinish\relax
\else
  \def\GDMfinish{\end{document}}
  \documentclass[12pt,a4paper]{article}
  \usepackage[margin=22mm]{geometry}
  \usepackage{amsmath,amssymb}
  \usepackage{graphicx}
  \usepackage{xepersian}
  \IfFontExistsTF{Amiri}{\settextfont{Amiri}}{\settextfont{DejaVu Sans}}
  \setlatintextfont{DejaVu Serif}
  \setlength{\parindent}{0pt}
  \setlength{\parskip}{6pt}
  \setlength{\emergencystretch}{3em}
  \begin{document}
\fi
\providecommand{\GDMm}[1]{\lr{\(#1\)}}
\providecommand{\GDMunclear}[1]{\textbf{[خوانش نامطمئن: #1]}}
\providecommand{\GDMpage}[1]{\clearpage\section*{صفحهٔ #1}}
\begin{center}\Large\bfseries دیابت بارداری ــ \lr{GDM}\end{center}
رونویسی تصاویر به ترتیب صفحات دست‌نویس؛ شاخه‌های فلش‌دار به صورت فهرست آمده‌اند. اعداد و عبارت‌های منبع حفظ شده‌اند و خوانش‌های نامطمئن مشخص شده‌اند.

\section*{صفحهٔ ۱}
\subsection*{هورمون‌های دیابتوژنیک}
\lr{GH}، \lr{CRH}، لاکتوژن جفتی (سوماتوماموتروپین کوریونیک)، پروژسترون، پرولاکتین.
\subsection*{تعریف‌ها}
\lr{DM}: کسی که قبل از هفتهٔ ۱۶ بارداری تشخیص داده شود، و یا قبل از بارداری \lr{DM} شناخته باشد.

۲--۳۸ درصد زنان باردار به \lr{GDM} مبتلا می‌شوند.
\subsection*{ریسک‌فاکتورهای \lr{GDM}}
\begin{itemize}
\item \lr{RF}های فردی: \lr{GDM} در بارداری قبلی، اختلال تحمل گلوکز، \GDMm{A_1C\geq5.7}، \lr{FBS} قبلی.
\item سابقهٔ خانوادگی \lr{DM} خصوصاً فامیل درجهٔ یک.
\item \GDMm{BMI\geq30} پیش از بارداری؛ افزایش وزن قابل توجه بین بارداری‌ها؛ افزایش وزن بیش از حد طی هفتهٔ ۱۸--۲۴ نخست بارداری.
\item بیماری‌های مرتبط با \lr{DM} مانند \lr{PCOS}.
\item سن مادر \GDMm{\geq35} سال.
\item سابقهٔ زایمان قبلی با \GDMm{BW\geq4\,kg}.
\end{itemize}
\subsection*{غربالگری زنان با \lr{Overt DM}}
از قبل \lr{DM} داشته یا \lr{DM} در اوایل بارداری ایجاد شده.
\begin{itemize}
\item چک \GDMm{HbA_1C} در نخستین ویزیت \lr{prenatal} برای همهٔ مادران باردار.
\item \GDMm{A_1C\geq6.5}: مادر مانند یک بیمار \lr{DM2} منیج شود.
\item \GDMm{A_1C<6.5}: در ۲۴--۲۸ هفتگی برای غربالگری \lr{GDM} تست ۲ مرحله‌ای انجام شود.
\item \GDMm{5.7\text{--}6.5}: با احتمال ۲۵\%؛ \GDMm{<5.7}: با احتمال \GDMm{<10\%} تشخیص \lr{GDM} می‌گیرند.
\end{itemize}
\subsection*{غربالگری \lr{GDM} در طی هفتهٔ ۲۴--۲۸}
همهٔ مادران باردار تست می‌شوند، غیر از آنهایی که تشخیص \lr{Overt DM} را دارند.

۲ رویکرد قابل انجام است:
\begin{enumerate}
\item تست یک‌مرحله‌ای: \lr{GTT} با ۷۵ گرم گلوکز ناشتا.
اگر ناشتا \GDMm{\geq92} یا یک ساعت \GDMm{\geq180} یا دو ساعت \GDMm{\geq153} باشد: \lr{GDM}.
\item تست دو‌مرحله‌ای: \lr{GTT} با ۵۰ گرم گلوکز در هر زمانی از روز.
\GDMm{\geq200}: \lr{GDM}؛ نتیجه بین \GDMm{135} (شامل) و \GDMm{200} (غیرشامل): مرحلهٔ دوم.
مرحلهٔ دوم: \lr{GTT} با ۱۰۰ گرم گلوکز ناشتا.
\end{enumerate}

\GDMpage{۲}
\subsection*{ادامهٔ تست دو‌مرحله‌ای}
اگر حداقل ۲ مورد از ۴ عدد مثبت شود: \lr{GDM}.
\begin{itemize}
\item ناشتا \GDMm{\geq95}؛ یک ساعت \GDMm{\geq180}؛ دو ساعت \GDMm{\geq155}؛ سه ساعت \GDMm{\geq140}.
\end{itemize}
تست‌هایی که نیاز به ناشتا بودن دارند: بیمار باید ۸--۱۰ ساعت ناشتا بوده باشد و نباید سیگار بکشد.

در صورتی که تست‌های غربالگری \lr{GDM} منفی شوند، توصیه می‌شود در دیگر بارداری تکرار شوند.

سطح خونی انسولین مادر ظرف چند روز پس از زایمان به سطح نرمال برمی‌گردد.

روز اول پس از زایمان قند خون بررسی شود؛ اگر نیاز بود انسولین داده شود.

پس از زایمان، مادر باید ظرف هفتهٔ ۴--۱۲ مجدد غربالگری شود و سپس سالانه تا ۳ سال: \GDMm{A_1C} و یا \lr{GTT} با ۷۵ گرم.

یک الگوریتم دیگر برای غربالگری \lr{Overt DM} در بارداری، طی هفتهٔ ۴--۱۴: افراد با \lr{RF} با \GDMm{A_1C} یا \lr{GTT} با ۵۰ گرم و کات‌آف \GDMm{\geq135} بررسی شوند.
\subsection*{درمان \lr{GDM}}
\begin{itemize}
\item ۸۵\% بیماران با \lr{GDM} می‌توانند با رژیم غذایی و تغییر سبک زندگی، نرموگلایسمیک شوند.
\item توصیه: قطع تمام نوشیدنی‌های قندی؛ به جای آن آب مصرف کنند.
\item محدود کردن میزان دریافت کربوهیدرات، خصوصاً قندهای ساده.
\item مصرف پروبیوتیک و رژیم‌های پر فیبر بی‌تأثیر است.
\item ۳ وعدهٔ غذایی کوچک تا متوسط و ۲--۴ نوبت اسنک.
\item رژیم کاهش وزن در مادر باردار توصیه نمی‌شود.
\item اگر \GDMm{BMI\geq25} و مادر \lr{GDM} باشد: کاهش وزن حدود ۳ کیلو قابل توصیه است.
\item فعالیت فیزیکی متوسط مانند پیاده‌روی ۳۰--۶۰ دقیقه روزانه در اکثر روزهای هفته به زنان \lr{GDM} توصیه می‌شود (حداقل هفتگی \lr{150 min})؛ مگر آنکه کنتراندیکاسیون مامایی وجود داشته باشد.
\end{itemize}
\subsection*{چک قند در بیمار \lr{GDM}}
روتین روزی ۴ نوبت توصیه می‌شود: ناشتا، و یک یا دو ساعت پس از شروع هر وعدهٔ غذایی؛ یک ساعت بهتر است.

در موارد \lr{mild GDM}، می‌توان یک بار در روز چک کرد و دفعات را کاهش داد؛ سپس پس از ۱ هفته. [عبارت منبع در همین‌جا پایان می‌یابد.]

تست‌ها اگر تنها ۵--۱۰ میلی‌گرم بالاتر از هدف باشند، شواهد رشد بیش از حد جنین وجود نداشته باشد، پس همین رژیم را ادامه می‌دهیم.

\GDMpage{۳}
در حال حاضر مطالعات کافی برای استفاده از \lr{CGM} (مانیتور مداوم) وجود ندارد.
\subsection*{تارگت گلوکز}
ناشتا: \GDMm{<95}؛ یک ساعت: \GDMm{<140}؛ دو ساعت: \GDMm{<120}.

عدد \GDMm{A_1C} به طور طبیعی در بارداری اندکی کاهش پیدا می‌کند.

چک مجدد \GDMm{A_1C} در بارداری \lr{GDM} جایگاهی ندارد؛ اما اگر انجام شد و با \GDMm{A_1C} اولیه هم‌خوانی نداشت، بررسی از نظر علت.

بررسی و مانیتورینگ کتونوری در \lr{GDM} جایگاهی ندارد.
\subsection*{فارماکوتراپی \lr{GDM}}
\begin{itemize}
\item هدف درمان: در اکثر تست‌ها قند خون کمتر از \lr{UL} باشد.
\item هیچ اپیزودی از هایپوگلایسمی ایجاد نشود.
\item کنترل شدید قند (قند متوسط کمتر از ۸۷) در \lr{GDM} کمک‌کننده نیست.
\end{itemize}
\textbf{اندیکاسیون‌ها:}
\begin{itemize}
\item بیماری با رژیم کنترل نشود: بیش از ۳۰\% قندها بالاتر از \lr{UL} باشند طی یک هفته.
\item شواهد \lr{overgrowth} جنین وجود داشته باشد: \lr{AC} بیشتر از صدک ۷۵ یا \lr{EFW} بیشتر از صدک ۹۰. [حاشیه: در مورد اوّلی مباحثه می‌شود.]
\item توصیه نمی‌شود که در زمان تشخیص \lr{GDM}، به صورت روتین متفورمین را شروع کنیم.
\end{itemize}
\subsection*{گزینه‌های فارماکوتراپی}
انسولین، متفورمین، گلیبوراید.

توصیه با انسولین به عنوان خط اول است.

برخی گایدلاین‌ها در مواردی که قند ناشتا طبیعی است و فقط قند پس از غذا بالاست، داروی خوراکی را نیز به عنوان خط اول می‌پذیرند.

درمان دو دارویی خوراکی در مقایسه با انسولین، ضعیف‌تر است و نباید استفاده گردد.

درمان با انسولین تنها است یا انسولین به اضافهٔ یک دارو.
\subsection*{دوز انسولین}
\begin{itemize}
\item معمولاً دوز مورد نیاز \GDMm{0.7\text{--}2\,U/kg} است (وزن فعلی بارداری).
\item حداقل ۴ \GDMunclear{نوبت متوالی؟} اندازه‌گیری (۴ نوبت روزانه یا بیشتر از آن ۹ نوبت).
\item پمپ انسولین در بارداری ممنوع است.
\item لیسپرو: بهترین \GDMunclear{دیتا؟} برای شروع انسولین نیست.
\end{itemize}

\GDMpage{۴}
\subsection*{شروع و تنظیم انسولین}
\begin{itemize}
\item شروع با تزریق ۱۰--۲۰ واحد انسولین متوسط‌الاثر مثل \lr{NPH} هنگام \lr{bedtime}، + ۶--۱۰ واحد انسولین سریع‌الاثر (لیسپرو، آسپارت، گلول) بلافاصله پیش از شروع غذا.
\item اگر قندهای تست پراندیال بالا باشند: افزایش ۱۰--۲۰ درصد \lr{Rapid-act}.
\item اگر فقط قند پس از شام بالا بود: ۴--۱۰ واحد \lr{Rapid} قبل از شام اضافه شود.
\item اگر فقط قند پس از ناهار بالا بود: ۴--۱۰ واحد \lr{Rapid} قبل از ناهار اضافه شود.
\item اگر قند صبح ناشتا و پس از ناهار بالا بود: بالا بردن دوز \lr{NPH}.
\item اگر قند ناشتا بالا بود و پس از اصلاح قند تست پراندیال‌ها: متوسط‌الاثر (\lr{NPH}) هنگام خواب با دوز \GDMm{0.2\,U/kg} اضافه می‌شود.
\end{itemize}
\subsection*{منیج هایپوگلایسمی در بارداری}
حد آن \GDMm{BS<60} است.
\begin{itemize}
\item معادل ۲۰--۳۰ گرم قند ساده به صورت خوراکی: ۵--۶ عدد قند یا نصف لیوان آب‌میوه.
\item ۱۵--۳۰ دقیقه بعد، مجدد چک قند.
\item ۲ نوبت هایپوگلایسمی: تغییر دوز انسولین.
\end{itemize}
انسولین‌های طولانی‌اثر در بارداری مطالعه نشده‌اند.

متفورمین و گلیبوراید هر دو از جفت رد می‌شوند، با غلظت‌های ۲ برابر و ۰٫۷ برابر نسبت به مادر؛ جنینی می‌شوند.

دوز متفورمین ۵۰۰ میلی‌گرم با غذا؛ به مرور، در صورت تحمل، اضافه شدن ۵۰۰ میلی‌گرم به همان وعده و یا بعدی است.

\lr{GDM} شانس ابتلا به \lr{DM2} را ۱۰ برابر بالا می‌برد.

ریسک کلی ابتلا به \lr{DM2} ظرف ۵ سال نخست بعد از زایمان ۹\% است و در تمام زندگی ۲۰--۵۰ درصد.
\subsection*{پیگیری‌های حین بارداری}
تقسیم‌بندی بیماران:
\begin{itemize}
\item \GDMm{A_1}: بدون نیاز به دارو؛ گلوکز به تارگت می‌رسد.
\item \GDMm{A_2}: با دارو.
\end{itemize}
\GDMm{A_1} بدون عارضه، مثل دیگران در بارداری: \lr{BPP} و سونوگرافی در هفتهٔ ۳۲--۳۴؛ ریسک افزایش‌یافتهٔ \lr{Stillbirth} ندارد.

\GDMpage{۵}
\subsection*{ادامهٔ پیگیری}
\GDMm{A_2} یا \GDMm{A_1} با کنترل قند \lr{Suboptimal}: \lr{NST}، \lr{AFI}، ۲ بار در هفته از هفتهٔ ۳۲.

در همهٔ \lr{GDM}ها یک سونو در هفتهٔ ۳۶--۳۹ لازم است.

ماکروزومی + پلی‌هیدرآمنیوس (\GDMm{EFW>4500}): برنامه‌ریزی برای \lr{C/S}.
\subsection*{زمان زایمان}
\begin{itemize}
\item در \GDMm{A_1} خوب: ۳۹ هفته تا ۴۱ هفته.
\item در \GDMm{A_1} یا \GDMm{A_2}: اینداکشن در هفتهٔ \GDMm{39^{+0}}.
\end{itemize}

\clearpage
\section*{تصاویر منبع}
\begin{center}
\textbf{تصویر ۱}\par
\IfFileExists{images/pic01.jpg}{\gdef\GDMOimagebase{images/}}%
  {\gdef\GDMOimagebase{Ob&Gyn/Oobstetrics/GDM/images/}}
\includegraphics[width=\textwidth,height=0.88\textheight,keepaspectratio]{\GDMOimagebase pic01.jpg}
\end{center}
\clearpage
\begin{center}
\textbf{تصویر ۲}\par
\includegraphics[width=\textwidth,height=0.88\textheight,keepaspectratio]{\GDMOimagebase pic02.jpg}
\end{center}
\clearpage
\begin{center}
\textbf{تصویر ۳}\par
\includegraphics[width=\textwidth,height=0.88\textheight,keepaspectratio]{\GDMOimagebase pic03.jpg}
\end{center}
\clearpage
\begin{center}
\textbf{تصویر ۴}\par
\includegraphics[width=\textwidth,height=0.88\textheight,keepaspectratio]{\GDMOimagebase pic04.jpg}
\end{center}
\clearpage
\begin{center}
\textbf{تصویر ۵}\par
\includegraphics[width=\textwidth,height=0.88\textheight,keepaspectratio]{\GDMOimagebase pic05.jpg}
\end{center}
\GDMfinish
